% sec-02-introduction.tex - Introduction (full)
\section{Introduction}

A Phase 1 oncology trial is, from a documentation standpoint, a relay of large
written artifacts. Before a single patient is dosed, an Investigational New Drug
(IND) dossier, an Investigator's Brochure, a clinical trial protocol, and an
informed consent form must be authored and approved; during the trial, safety
narratives, amendments, and cohort-review packages must be produced on short
clocks; and after the trial, a clinical study report and downstream regulatory
modules must be assembled. Each of these documents can sit on the critical path of
the next decision. This paper (paper v1.0, within repository v4.2.0) shows that a
single repository based large language model (LLM) prompt, which first writes and
then executes a schedule of sub-prompts, can generate every one of these documents
through a mermaid to draft to full to final pipeline, and can do so in a way that
is monitorable and verifiable in real time.

\subsection{Patients Currently Enrolled in Trial Cannot Wait for Slowdowns}

Pancreatic ductal adenocarcinoma (PDAC) has a five-year overall survival below 13
percent and is among the most lethal solid tumors, so a patient already enrolled in
a Phase 1 PDAC study has little survival margin to spare on administrative delay
\cite{kawchak_2025_15735068, natplatform}. The clinical and operational time of a
trial (recruitment, treatment, and waiting for progression-free and overall
survival to mature) cannot be compressed by writing faster, and neither can the
fixed regulatory review clocks. What can be compressed is the administrative and
preparation time spent authoring the large documents between steps. Faster
paperwork therefore means faster treatment: an internally consistent package can
reach the IRB and the FDA sooner, sites can activate sooner, and the next cohort can
open sooner. In this work, a new iteration of any large document (a fresh draft, a
full revision, or a final polish) completes in one to four days rather than the
weeks to months a sequential human authoring cycle typically consumes. Figure 1
shows the single-prompt pipeline that produces these iterations, and Figure 2 shows
the one-to-four-day cadence against the traditional baseline.

\begin{mermaidfig}
\node[mmgoal] (mp) at (0,0) {Master prompt\\prompts/prompt-paper.md\\paper v1.0, repo v4.2.0};
\node[mmproc] (a) at (0,-1.95) {Process A\\generate sub-prompts 1-4};
\node[mmin] (s1) at (0,-3.9) {Stage 1 mermaid\\24 colored figures};
\node[mmin] (s2) at (0,-5.85) {Stage 2 draft-paper\\scaffold + bracketed instr.};
\node[mmaccent] (s3) at (0,-7.8) {Stage 3 full-paper\\prose + TikZ + tables};
\node[mmgoal] (s4) at (0,-9.75) {Stage 4 final-paper\\polished + Overleaf zip};
\node[mmproc] (rel) at (0,-11.7) {Repository update\\README + releases + CHANGELOG v4.2.0};
\node[mmctx] (gh) at (6.0,-5.85) {GitHub\\real-time\\auto-commit / auto-PR\\monitor, no intervention};
\draw[mmedge] (mp)--(a);
\draw[mmedge] (a)--(s1);
\draw[mmedge] (s1)--(s2);
\draw[mmedge] (s2)--(s3);
\draw[mmedge] (s3)--(s4);
\draw[mmedge] (s4)--(rel);
\draw[mmedged] (s1.east) to[out=0,in=150] (gh.north west);
\draw[mmedged] (s2.east) -- (gh.west);
\draw[mmedged] (s3.east) to[out=0,in=210] (gh.south west);
\draw[mmedged] (s4.east) to[out=0,in=270] (gh.south);
\draw[mmedged] (gh.north) to[out=90,in=20] (s1.east);
\end{mermaidfig}
\figcaption{Figure 1. The single-prompt mermaid to draft to full to final build
pipeline. One master prompt drives Process A (generate the four sub-prompts) and
then Process B (execute them in order). Dashed edges show that every stage
auto-commits and opens a pull request on GitHub in real time, and the curved
feedback edge shows author review returning into the build, all without manual
intervention.}

\begin{mermaidfig}
\node[mmdark] (t0) at (0,0) {Traditional baseline\\weeks to months\\per large-document iteration};
\node[mmctx] (t1) at (0,-1.8) {sequential drafting};
\node[mmctx] (t2) at (0,-3.3) {manual table assembly};
\node[mmctx] (t3) at (0,-4.8) {serial review rounds};
\node[mmaccent] (n0) at (7.2,0) {Single-prompt method\\1 to 4 days\\per new iteration};
\node[mmin] (n1) at (7.2,-1.8) {draft scaffold (day 1)};
\node[mmin] (n2) at (7.2,-3.3) {full version (days 2-3)};
\node[mmin] (n3) at (7.2,-4.8) {final polish (day 4)};
\node[mmgoal] (g) at (3.6,-6.6) {Administrative / prep\\time compressed};
\draw[mmedge] (t0)--(t1); \draw[mmedge] (t1)--(t2); \draw[mmedge] (t2)--(t3);
\draw[mmedge] (n0)--(n1); \draw[mmedge] (n1)--(n2); \draw[mmedge] (n2)--(n3);
\draw[mmedge] (n3) to[out=-90,in=0] (g.east);
\draw[mmedged] (t3) to[out=-90,in=180] (g.west);
\end{mermaidfig}
\figcaption{Figure 2. New document iterations in one to four days. The traditional
baseline (left) spends weeks to months per large-document iteration in sequential
drafting, manual table assembly, and serial review. The single-prompt method
(right) produces a draft, a full version, and a final polish in one to four days,
compressing the administrative and preparation bucket only.}

\subsection{Effective Verifications Prove Utility}

Speed is worthless if it cannot be trusted, so the method is built to be verified
rather than asserted. Five checks make the build self-evidently grounded. First, the
prose is grounded in 24 colored Mermaid figures that carry the real quantitative
content (document counts, reporting windows, and review clocks) rather than
decoration. Second, each figure matches the surrounding paper context and is cited
where it is discussed. Third, every generated file is committed to GitHub in real
time, so the author can view each artifact as it appears on the branch. Fourth, the
paper's repository and directory names match the repository names exactly, so a
reader can navigate from a sentence to the directory it describes. Fifth, the
paper's file names match the repository file names, so a cited
\texttt{sections/sec-04-results.tex} is the file on disk. These five checks are
collected in the Results.

\subsection{Patient Lives Extended: Probable Benefit Exceeds Probable Risk}

The motivation is ultimately clinical. The probable benefits of the method are
faster paperwork, faster treatment, a compressed administrative and preparation
bucket, and more figures available for human oversight. The probable risks are the
familiar LLM failure modes: an occasional generation error, a formatting artifact,
or an unverified claim. Each risk is reduced to a low residual by a process control:
human-in-the-loop review at every commit, Mermaid grounding, paper-to-repository
name matching, real-time GitHub monitorability, and a senior-author proofreading
pass. Figure 3 weighs the two sides. For a population whose disease will not wait,
withholding a method that safely compresses preparation time carries its own cost,
and the balance favors using it.

\begin{mermaidfig}
\node[mmin] (r1) at (0,0) {Possible LLM error\\(server / oversized input)};
\node[mmin] (r2) at (0,-1.7) {Formatting artifacts};
\node[mmin] (r3) at (0,-3.4) {Hallucinated context};
\node[mmproc] (m1) at (4.7,0) {Human-in-the-loop\\review per commit};
\node[mmproc] (m2) at (4.7,-1.7) {Mermaid grounding\\and name-matching};
\node[mmproc] (m3) at (4.7,-3.4) {Real-time GitHub\\and author proofread};
\node[mmdec] (d) at (9.2,-1.7) {Benefit vs\\risk};
\node[mmgoal] (g) at (9.2,-4.2) {Benefit $>$ risk\\for patients who\\cannot wait};
\node[mmaccent] (b) at (9.2,0.6) {Faster paperwork,\\faster treatment};
\draw[mmedge] (r1)--(m1); \draw[mmedge] (r2)--(m2); \draw[mmedge] (r3)--(m3);
\draw[mmedge] (m1) to[out=0,in=120] (d.north west);
\draw[mmedge] (m2)--(d);
\draw[mmedge] (m3) to[out=0,in=240] (d.south west);
\draw[mmedge] (b)--(d);
\draw[mmedge] (d)--(g);
\end{mermaidfig}
\figcaption{Figure 3. Probable benefit exceeds probable risk. Each probable risk
(left) is reduced to a low residual by a process control (center); the decision
node (right) integrates the mitigated risks with the weighted benefits and resolves
in favor of using the method for enrolled patients who cannot wait.}

\subsection{Scope, Contributions, and Relationship to the Adoption Guide}

This paper builds directly on the \emph{Oncology Trial PI LLM Adoption Guide}
\cite{adoptionguide}, which it uses both as the paper template and as the practical
hands-on guide whose Repository Setup, Prompt Proficiency, Project Proficiency, and
Figure and Table guidance it turns into a worked Phase 1 PDAC document-generation
study. The contributions are fourfold: (1) a single-prompt mermaid to draft to full
to final pipeline that produces a complete scientific paper and, by the same
mechanism, the large Phase 1 trial documents; (2) explicit coverage of the six
document targets where faster authoring yields the greatest schedule value; (3)
twenty-four new, professionally colored figures that reproduce identically from
Python Mermaid to LaTeX TikZ; and (4) a real-time, monitorable, name-matched build.
The demonstration is anchored to a real-world program: the on-premises LLM directed
robotic Whipple with perioperative daraxonrasib (RMC-6236), whose Phase 1
\cite{phase1protocol} and Phase 2 \cite{phase2protocol} protocols and whose
60-second simulation \cite{kawchak_2026_20196639} were themselves single-prompt
multi-file builds, and to the National Platform for Physical AI Oncology Trials as
the governing framework \cite{natplatform}.
